%% outlier_detection.tex -- the original Tensor Omics analysis: which genes have
%% diverged from their family, and in which direction.

\begin{recipe}[label=rec:outlier-detection]{Detecting Divergent Genes: Outliers and Shift Vectors}

\recipepurpose{Find the genes whose expression has moved away from the rest of
  their gene family further than the family's own spread explains, and describe
  \emph{how} each one moved. This is the original Tensor Omics workflow: the
  distances answer whether a copy diverged, the shift vectors answer in which
  direction.}

\nomaintainer

\begin{requiredinputs}
  \item Normalized expression vectors, one gene per column, from the
        normalization sub-recipe.
  \item A gene-to-family mapping, with \code{0} for genes belonging to no
        family. Genes outside your families may stay in the matrix.
  \item Family centroids, from the centroid sub-recipe --- and the decision
        made there about which members define them.
\end{requiredinputs}

\begin{prerequisites}
  \dependson{sub:normalization}
  \dependson{sub:family-centroids}
\end{prerequisites}

\begin{workflow}
  \step{Measure each gene against its own family's centroid with
        \code{distance\_to\_centroid}. Genes with no family receive
        \code{-1}, a sentinel rather than a distance.}
  \step{Call \code{detect\_outliers}. It does \emph{not} compare distances to a
        fixed cutoff: it divides each distance by its family's own spread ---
        the Relative Distance Index (RDI) --- and flags the top
        $1-\optn{percentile}$ of those. Scaling per family is what puts a
        broadly spread family and a tight one on the same footing.}
  \step{Compute the shift vector field with
        \code{compute\_shift\_vector\_field}. For each gene it returns two
        vectors, the family centroid and then the shift
        $\vec{s} = \vec{p} - \vec{o}$ from that centroid to the gene.}
  \step{Report the flagged genes with their shift vectors. Use the
        \emph{raw} distances and shifts for everything downstream --- the RDI
        exists to choose which genes are worth looking at, not to describe how
        far they moved.}
\end{workflow}

\begin{note}
  The family scaling is not simply each family's own standard deviation. A
  family with few members cannot give a usable one, so Tensor Omics fits the
  spread of a family against its mean distance across \emph{all} families with
  LOESS, and reads each family's scale off that curve. A small family is
  therefore scaled by the trend its neighbours establish instead of by its own
  noise.
\end{note}

\examplescript{python/examples/outlier_detection.py}{%
  Runs the whole chain --- centroids, distances, outliers, shift vectors ---
  and writes a per-gene table plus the shift vectors of the flagged genes.
  Run unchanged from the repository root it takes about two seconds on the
  bundled data and writes \file{results/outliers.tsv}.}

\begin{codefragment}[language=Python,caption={Fragments to edit: the threshold, and what you keep}]
from tensor_omics import (distance_to_centroid, detect_outliers,
                          compute_shift_vector_field)

# --- how far is each gene from its own family's centroid? --
distances = distance_to_centroid(
    expr, centroids, gene_to_family)   # -1 where no family

# --- which of those distances are extreme for their family? -
result = detect_outliers(
    n_families, distances, gene_to_family,
    percentile=0.95)                   # FRACTION: top 5%
is_outlier = result["is_outlier"]
quantile   = result["quantile"]        # effect size, NOT a p-value

# --- in which direction did each gene move? ----------------
field = compute_shift_vector_field(
    expr, centroids, gene_to_family)   # (n_axes, 2, n_genes)
centroid_of = field[:, 0, :]           # the family's centroid
shift       = field[:, 1, :]           # p - o, the divergence
\end{codefragment}

\begin{parameters}
  \param{percentile}{Where the threshold on the RDI sits; the genes above it
    are flagged}{A \emph{fraction} in $[0,1]$, not a percentage:
    \code{0.95} keeps the top 5\%. Lower it to screen more broadly, raise it to
    report only the strongest divergence}
  \param{n\_families}{How many families the mapping refers to}{Must cover the
    largest index in \code{gene\_to\_fam}, and must be the same number the
    centroids were computed with}
  \param{gene\_to\_fam}{Which family each gene is measured against}{\code{0}
    for a gene belonging to no family. Such genes pass through the whole chain
    without being flagged and without affecting the threshold}
\end{parameters}

\begin{expectedresults}
  With the default \code{percentile=0.95} the flagged fraction is essentially
  5\% of the genes that have a family --- on the bundled data, 512 of 10\,234,
  exactly 5.00\%, with 75\,594 genes belonging to no family carried along and
  none of them flagged. Distances to the own centroid have a median of 1.73 and
  a maximum of 2.96 there. A flagged fraction far from
  $1-\optn{percentile}$ means genes are reaching the test that should not:
  check for families with a zero centroid first.
\end{expectedresults}

\begin{pitfall}[Pitfall: reading percentile as a percentage]
  \optn{percentile} is a fraction in $[0,1]$. Passing \code{95} does not ask
  for the top 5\% --- it is out of range and the call fails. Passing
  \code{0.05} silently asks for the top 95\%, flagging almost everything.
\end{pitfall}

\begin{pitfall}[Pitfall: treating the quantile as a p-value]
  The \code{quantile} output states how extreme a gene's distance is
  \emph{relative to the distances actually observed} --- an empirical
  upper-tail effect size. It is not a null-hypothesis test: nothing here models
  what the distances would look like in the absence of divergence, so a
  \code{quantile} of 0.001 is not a $p$ of 0.001 and must not be corrected for
  multiple testing as though it were.
\end{pitfall}

\begin{pitfall}[Pitfall: a family whose spread cannot be estimated]
  A family with a single member \emph{is} its own centroid: the distance is
  zero, there is no intra-family spread to estimate, its scaling factor is
  zero, and the gene can never be flagged. That is the right answer --- one
  copy cannot diverge from itself --- but it is silent, so a family reduced to
  one member by filtering leaves the analysis without a word. Check family
  sizes rather than trusting an empty result. And where too few families remain
  to fit the trend at all, every family falls back to a single global spread,
  losing exactly the between-family comparability the RDI exists to provide.
\end{pitfall}

\begin{pitfall}[Pitfall: testing genes against a reference that contains them]
  If the centroids came from \code{group\_centroid\_all}, every candidate is
  part of the reference it is measured against, which pulls the centroid
  towards the divergent copies and shrinks exactly the distances this recipe
  looks for. Prefer orthologs-only centroids when the question is whether a
  copy diverged --- see \seerecipe{sub:family-centroids}.
\end{pitfall}

\begin{pitfall}[Pitfall: averaging the distance sentinel]
  \code{distance\_to\_centroid} returns \code{-1} for a gene with no family.
  That is a marker, not a short distance: summarising the distance column
  without excluding it drags every statistic down. Select on the family
  mapping, not on the distances.
\end{pitfall}

\begin{interpretation}
  A flagged gene is one that sits further from its family's consensus than that
  family's own spread makes ordinary. It is a statement about \emph{relative}
  divergence: the same absolute distance is unremarkable in a widely spread
  family and striking in a tight one, which is the whole point of scaling per
  family.

  The distance says only \emph{that} a gene moved. The shift vector says
  \emph{where} to, and that is the interpretable part: its large components
  name the tissues that carry the divergence. A shift concentrated on one axis
  is a copy that gained or lost expression in that one tissue; a shift spread
  evenly across axes is a change in overall level rather than in pattern. Read
  the shift vectors of the flagged genes before drawing conclusions --- and
  carry the raw shifts, not the RDI, into any later geometric analysis.

  Finally, the flagged set is a \emph{ranking cut}, not a discovery set with an
  error rate. At \code{percentile=0.95} you get the top 5\% by construction,
  whether or not anything in your data has genuinely diverged.
\end{interpretation}

\begin{references}
  \reference{Cleveland (1979). Robust locally weighted regression and
    smoothing scatterplots. \emph{Journal of the American Statistical
    Association} 74(368):829--836.}
\end{references}

\end{recipe}
